\documentclass[preprint,10pt,5p]{elsarticle}

\usepackage{amssymb}
\usepackage{lineno}
\usepackage{amsmath}
\usepackage{subcaption}
\usepackage{hhline}
\usepackage{hyperref}
\usepackage{booktabs}
\usepackage[utf8]{inputenc}
\usepackage[T1]{fontenc}
\usepackage[brazil]{babel}

\journal{Universidade Católica Dom Bosco}

\begin{document}

\begin{frontmatter}

\title{Nexus: Sistema de Análise de Dados Educacionais com Suporte de Inteligência Artificial para Monitoramento da ODS~4 no Mato Grosso do Sul}

\address[label1]{Universidade Católica Dom Bosco, Campo Grande, Brasil}

\author[label1]{Abner Matheus Moraes Saconi}
\ead{abnermatheusmoraes@gmail.com}

\begin{abstract}
O acompanhamento quantitativo dos indicadores educacionais da rede pública estadual do
Mato Grosso do Sul constitui requisito indispensável para a verificação das metas do
Objetivo de Desenvolvimento Sustentável número~4 (ODS~4) da Agenda 2030 das Nações Unidas.
Os dados disponibilizados pela Secretaria de Estado de Educação encontram-se fragmentados em
arquivos CSV anuais, inviabilizando análises históricas sem processamento especializado.
Este trabalho descreve o desenvolvimento do Nexus, sistema web que consolida a importação,
persistência e análise desses registros por meio de um módulo denominado DAC
(\textit{Dados Acadêmicos das Escolas}), estruturado em hierarquia orientada a objetos
(\texttt{DadoAnual}, \texttt{Escola}, \texttt{Municipio}, \texttt{SerieHistorica}) e
integrado a banco de dados relacional com cobertura de 2018 a 2024.
A base resultante totaliza mais de 8.700 registros únicos distribuídos por 1.200 unidades
escolares em 79 municípios. A análise dos dados importados evidencia redução acumulada de
32,3\% no total de matrículas ao longo do período, recuperação das taxas de aprovação entre
2018 e 2023 — de 81,3\% para 88,7\% — e elevação expressiva da taxa de reprovação em 2024
(22,8\%), cujas causas merecem investigação aprofundada. Um módulo de chat analítico,
alimentado por modelo de linguagem de grande escala com contexto construído dinamicamente
a partir dos filtros ativos, possibilita a formulação de consultas em linguagem natural
diretamente sobre esses dados, sem exigir conhecimento prévio em análise quantitativa.
\end{abstract}

\begin{keyword}
educação de qualidade \sep ODS~4 \sep análise de dados educacionais \sep inteligência artificial \sep
Mato Grosso do Sul \sep sistema web
\end{keyword}

\end{frontmatter}


%-----------------------------------------------------------------------
\section{Introdução}
\label{intro}
%-----------------------------------------------------------------------

Entre os 17 Objetivos de Desenvolvimento Sustentável estabelecidos pela Agenda 2030 das Nações
Unidas, o ODS~4 — Educação de Qualidade — ocupa posição central, por condicionar o alcance de
vários outros objetivos ao nível de instrução da população~\cite{onu2015}. Seu monitoramento
efetivo exige a sistematização periódica de indicadores como taxas de aprovação, reprovação e
abandono escolar, cuja interpretação isolada, sem série histórica, oferece subsídios limitados
para a formulação de políticas públicas.

No Mato Grosso do Sul, a Secretaria de Estado de Educação (SED-MS) publica anualmente os
dados de matrículas por unidade escolar em formato CSV aberto~\cite{seduc_ms}. Cada arquivo
reúne dezenas de milhares de registros com situação final dos alunos por escola, contudo sem
encadeamento entre exercícios. A ausência de uma ferramenta que consolide esses arquivos em
série histórica única impede comparações interanuais e análises de tendência — operações
essenciais para avaliar o progresso em direção às metas do ODS~4.

Diante desse cenário, o presente trabalho relata o desenvolvimento do \textbf{Nexus},
sistema de inteligência artificial acadêmica concebido no âmbito do Desafio de Articulação
de Competências (DAC) do 4º semestre dos cursos de Ciências da Computação e Tecnologia em
Análise e Desenvolvimento de Sistemas da Universidade Católica Dom Bosco (UCDB). O módulo
central do sistema — denominado DAC — provê quatro capacidades integradas: (i)~importação
automatizada e fusão dos CSVs anuais em banco relacional único com deduplicação idempotente;
(ii)~dashboard interativo com filtros dinâmicos por município e intervalo de anos;
(iii)~hierarquia orientada a objetos para agregação e cômputo de indicadores históricos; e
(iv)~chat analítico baseado em LLM cujo contexto é construído dinamicamente a partir dos
filtros ativos pelo usuário.

O artigo está organizado da seguinte forma: a Seção~\ref{works} situa o trabalho na
literatura correlata; a Seção~\ref{methods} detalha os dados, a arquitetura e o ambiente
experimental; a Seção~\ref{results} apresenta e discute os resultados obtidos;
a Seção~\ref{discussion} analisa limitações e contribuições; e a Seção~\ref{conclusion}
sintetiza as conclusões e aponta direções de pesquisa futura.


%-----------------------------------------------------------------------
\section{Trabalhos Relacionados}
\label{works}
%-----------------------------------------------------------------------

A mineração de dados educacionais (\textit{Educational Data Mining} — EDM) constitui campo
consolidado, conforme atestam \citet{romero2010} em revisão abrangente que cataloga técnicas
de agrupamento, classificação e detecção de padrões de risco aplicadas a registros de
ambientes virtuais de aprendizagem. A distinção fundamental em relação ao presente trabalho
reside na natureza dos dados: enquanto o corpus da EDM é tipicamente transacional e gerado
por plataformas LMS, o módulo DAC opera sobre registros administrativos anuais de matrículas,
com granularidade por escola e periodicidade determinada por calendário escolar, o que implica
volume menor de registros porém abrangência territorial integral.

No plano institucional, o Instituto Nacional de Estudos e Pesquisas Educacionais Anísio
Teixeira produziu relatório específico sobre o ODS~4 no Brasil~\cite{inep2022}, baseado nas
bases de dados do Censo Escolar. As análises ali apresentadas, todavia, são conduzidas por
equipes técnicas especializadas com ferramentas proprietárias e disponibilizadas apenas como
relatórios estáticos, sem qualquer interface interativa ou possibilidade de recorte por
município.

A integração de LLMs como camada de consulta sobre bases estruturadas tem sido investigada
por \citet{liu2023}, que demonstram ser possível atingir precisão elevada em perguntas
analíticas quando o modelo recebe, no contexto, um resumo suficientemente completo do
subconjunto de dados relevante. O Nexus implementa essa estratégia de forma pragmática: o
método \texttt{resumo\_para\_llm()} serializa os agregados calculados pela hierarquia OOP e
os injeta no \textit{prompt} do sistema antes de cada resposta, eliminando a necessidade de
\textit{fine-tuning} e reduzindo o risco de alucinações factuais.

Sistemas como o QEdu~\cite{qedu2023} e o Painel de Monitoramento do INEP oferecem
visualizações de indicadores nacionais, sem suporte a análises históricas municipais
aprofundadas nem integração com assistentes conversacionais. O Nexus preenche essa lacuna
ao concentrar-se nos dados estaduais do Mato Grosso do Sul e ao combinar análise visual
com linguagem natural num único ambiente operacional.


%-----------------------------------------------------------------------
\section{Materiais e Métodos}
\label{methods}
%-----------------------------------------------------------------------

\subsection{Conjunto de Dados}

Os dados processados pelo módulo DAC provêm dos relatórios anuais de matrículas por unidade
escolar publicados pela SED-MS no portal estadual de dados abertos~\cite{seduc_ms}.
Cada arquivo CSV reúne registros individuais por escola contendo: código e nome da unidade
escolar, município, ano de referência, total de matrículas inicial e pós-censo, e situação
dos alunos ao término do ano letivo — aprovados, reprovados, em abandono, transferidos,
cancelados, ainda cursando e outras situações.

Foram importados arquivos referentes aos anos de 2018 a 2024, perfazendo sete exercícios com
aproximadamente 3.700 registros cada. O arquivo de 2026 foi excluído da análise por não conter
resultados finais, e o ano de 2025 permanecia inédito à época da realização deste trabalho.
Após a fusão e deduplicação, o banco de dados resultante totaliza 8.742 registros únicos,
indexados pelo par \texttt{(escola\_id, ano)}, cobrindo 1.247 unidades escolares distribuídas
em 79 municípios do estado.

\subsection{Arquitetura do Sistema}

O Nexus adota arquitetura em quatro camadas com responsabilidades bem delimitadas.

\textbf{Camada de dados.} Implementada em Python com FastAPI~\cite{fastapi} e
SQLAlchemy~\cite{sqlalchemy} sobre banco SQLite. O esquema relacional é composto por duas
tabelas: \texttt{dac\_escolas}, com restrição de unicidade sobre \texttt{(nome, municipio)},
e \texttt{dac\_dados\_escolares}, com restrição \texttt{UNIQUE(escola\_id, ano)} que garante
idempotência nas reimportações. O pipeline detecta automaticamente a codificação do arquivo
(UTF-8, Latin-1 ou CP1252) e o delimitador CSV, normaliza os cabeçalhos por remoção de
acentos e caracteres especiais, e executa inserções em lote via \textit{bulk insert} com
deduplicação em memória.

\textbf{Camada de domínio.} Uma hierarquia de quatro classes agrega os registros em memória
para cômputo de indicadores. \texttt{DadoAnual} encapsula os contadores brutos de um ano
letivo em uma escola e calcula as taxas de aprovação, reprovação e abandono sobre a base de
alunos com resultado definido. \texttt{Escola} mantém o histórico anual de uma unidade
escolar. \texttt{Municipio} agrega as escolas, computa totais e taxas municipais e expõe o
método \texttt{taxa\_reprovacao(ano)}. \texttt{SerieHistorica} representa o estado completo
da base e produz, via \texttt{resumo\_para\_llm()}, o resumo textual estruturado injetado
como contexto no LLM.

\textbf{Camada de apresentação.} Desenvolvida em React com TypeScript e biblioteca Recharts.
O dashboard apresenta KPIs calculados dos totais brutos (matrículas, aprovação, reprovação e
abandono), gráficos de linha para evolução histórica e gráfico de barras para comparativo
anual, além de seletor dinâmico de colunas. Os filtros por município e por intervalo de anos
atualizam todos os componentes sem recarregamento de página.

\textbf{Camada de inteligência artificial.} O chat analítico integra o modelo
Hermes-3-Llama-3.1-8B quantizado em Q4\_K\_M, servido via llama.cpp~\cite{llamacpp}.
O contexto injetado no \textit{prompt} é gerado a partir dos filtros ativos, permitindo
que o modelo opere exclusivamente sobre o subconjunto de dados selecionado pelo usuário.
As respostas são transmitidas ao cliente em tempo real via \textit{Server-Sent Events} (SSE).

\subsection{Ambiente Experimental}

O sistema foi implantado por meio de Docker Compose em quatro serviços isolados: servidor
LLM, servidor de embeddings, \textit{backend} FastAPI e \textit{frontend} React/Vite.
O serviço LLM operou em uma GPU NVIDIA GeForce RTX~3070 (8~GB de VRAM), suficiente para
a quantização Q4\_K\_M do modelo de 8 bilhões de parâmetros. Queries SQLAlchemy de grande
volume utilizam \texttt{yield\_per(1000)} para evitar alocação integral do resultado em
memória, viabilizando a construção do contexto do LLM sem restrições de escala.

A validação funcional consistiu na importação sequencial dos sete arquivos CSV, verificação
da idempotência mediante reimportação e teste do chat analítico com perguntas sobre tendências
históricas, comparações intermunicipais e projeções qualitativas, contrastando as respostas
do LLM com os valores calculados diretamente pela hierarquia OOP.


%-----------------------------------------------------------------------
\section{Resultados}
\label{results}
%-----------------------------------------------------------------------

A importação dos sete arquivos CSV foi concluída em tempo inferior a cinco segundos, com
todos os registros disponíveis imediatamente no dashboard. O mecanismo de deduplicação
mostrou-se efetivo: a reimportação completa dos arquivos não alterou nenhum registro já
persistido, confirmando a idempotência do pipeline.

A Tabela~\ref{table:dados} consolida as principais métricas estaduais extraídas da base
importada para o período 2018--2024.

\begin{table*}[t]
\centering
\caption{Dados educacionais da rede pública estadual do Mato Grosso do Sul importados no Nexus
(2018--2024). Taxas calculadas sobre alunos com resultado definido (aprovados + reprovados + abandono).
Fonte: SED-MS~\cite{seduc_ms}.}
\label{table:dados}
\small
\begin{tabular}{|c|r|r|r|r|}
\hline
\textbf{Ano} & \textbf{Matrículas} & \textbf{Aprovação} & \textbf{Reprovação} & \textbf{Abandono} \\
\hline\hline
2018 & 326.367 & 81,3\% & 14,6\% & 4,0\% \\
\hline
2019 & 293.195 & 86,7\% & 13,2\% & 0,1\% \\
\hline
2020 & 242.028 & 89,4\% & 10,6\% & $-$    \\
\hline
2021 & 248.780 & 88,4\% & 11,6\% & $-$    \\
\hline
2022 & 247.646 & 88,0\% & 12,0\% & $-$    \\
\hline
2023 & 237.574 & 88,7\% & 11,3\% & $-$    \\
\hline
2024 & 220.926 & 77,2\% & 22,8\% & $-$    \\
\hline
\end{tabular}

\vspace{2pt}
\noindent\footnotesize{($-$) Coluna de abandono não reportada no arquivo oficial a partir de 2020.
O ano de 2024 apresenta taxa de reprovação elevada, possivelmente relacionada à atualização
dos critérios de avaliação adotados pela rede estadual naquele período.}
\end{table*}

Três tendências merecem destaque a partir da leitura da tabela. Primeiro, verifica-se
redução acumulada de 32,3\% no total de matrículas entre 2018 e 2024 — de 326.367 para
220.926 —, fenômeno consistente com a retração demográfica observada em estados de médio
porte e com a migração de alunos para redes municipais ou federais. Segundo, as taxas de
aprovação e reprovação exibiram trajetória de melhora progressiva entre 2018 e 2020, período
em que a aprovação saltou de 81,3\% para 89,4\% e a reprovação recuou de 14,6\% para 10,6\%.
A inflexão observada em 2020 coincide com a suspensão das aulas presenciais e a adoção de
critérios de avaliação excepcionais durante a pandemia de Covid-19, o que pode ter contribuído
artificialmente para os resultados daquele ano. Terceiro, e de forma mais preocupante, o
exercício de 2024 registra taxa de reprovação de 22,8\% — patamar superior até mesmo ao de
2018, antes do período de melhora —, o que demanda investigação sobre eventuais mudanças nos
critérios de avaliação adotados pela SED-MS.

No que concerne ao módulo de chat analítico, as consultas formuladas em linguagem natural
produziram respostas coerentes com os valores tabulados, sem alucinações factuais detectadas
nas verificações realizadas. O primeiro \textit{token} de cada resposta foi exibido com
latência média inferior a dois segundos, viabilizando interação fluida mesmo com o contexto
completo da série histórica injetado no \textit{prompt}.


%-----------------------------------------------------------------------
\section{Discussão}
\label{discussion}
%-----------------------------------------------------------------------

A análise dos dados obtidos com o Nexus revela que a série histórica 2018--2024 da rede
pública estadual do Mato Grosso do Sul não segue trajetória linear nem uniforme. A melhora
consistente das taxas entre 2018 e 2023 poderia, isoladamente, sugerir avanço estrutural no
desempenho escolar; contudo, a reversão abrupta de 2024 — com reprovação saltando de 11,3\%
para 22,8\% em um único exercício — impede qualquer conclusão otimista sem escrutínio
adicional das condições metodológicas de coleta. Cabe investigar, em estudos futuros, se essa
variação reflete mudança real de desempenho ou reajuste nos critérios de classificação de
alunos adotados pela secretaria.

Do ponto de vista técnico, a estratégia de deduplicação baseada em restrições \texttt{UNIQUE}
no nível do banco — em vez de verificações em camada de aplicação — mostrou-se tanto mais
eficiente quanto mais robusta, pois delega ao mecanismo relacional a garantia de consistência.
O uso de \texttt{yield\_per(1000)} nas queries de construção de contexto permitiu processar
a totalidade dos 8.742 registros sem pressão sobre a memória do contêiner, confirmando a
adequação da abordagem para volumes maiores.

Uma limitação relevante é a dependência de mapeamento estático entre os cabeçalhos dos CSVs
e os campos do esquema relacional. A normalização implementada — que remove acentos,
underscores e espaços — mitiga variações ortográficas menores, mas não acomoda renomeações
estruturais de colunas. Outra limitação diz respeito ao escopo estadual: a cobertura atual
restringe-se à rede pública estadual do Mato Grosso do Sul, excluindo as redes municipais
e federal. A ausência da coluna de abandono nos arquivos a partir de 2020 constitui lacuna
nos dados de origem que o sistema não pode suprir sem informação complementar da SED-MS.

Por fim, o LLM empregado — 8 bilhões de parâmetros em quantização de 4 bits — demonstra
capacidade adequada para análises descritivas e comparativas, porém não deve ser utilizado
para inferências causais ou projeções quantitativas sem validação explícita dos resultados.
A arquitetura em camadas adotada permite a substituição do modelo por versão de maior
capacidade sem alterações na cadeia de coleta e agregação de dados.


%-----------------------------------------------------------------------
\section{Conclusão}
\label{conclusion}
%-----------------------------------------------------------------------

Este trabalho apresentou o Nexus, sistema web de análise de dados educacionais com módulo
de inteligência artificial integrado, desenvolvido para consolidar e tornar analiticamente
acessíveis os registros anuais de matrículas da rede pública estadual do Mato Grosso do Sul.
A partir da importação de sete arquivos CSV anuais, o sistema construiu uma base relacional
com 8.742 registros únicos cobrindo 1.247 escolas em 79 municípios — base sobre a qual foram
identificadas tendências relevantes: contração de 32,3\% no volume de matrículas ao longo de
seis anos, recuperação das taxas de aprovação entre 2018 e 2023, e elevação expressiva da
reprovação em 2024, ainda sem explicação consolidada na literatura.

A combinação de persistência relacional com deduplicação automática, hierarquia orientada a
objetos para agregação de indicadores e chat analítico com contexto construído dinamicamente
constitui contribuição metodológica que viabiliza o monitoramento contínuo da ODS~4 em nível
estadual sem exigir expertise técnica do usuário final. O código-fonte, a estrutura de dados
e o ambiente Docker compõem um artefato reproduzível que pode ser adaptado a outros estados
ou ampliado com dados das redes municipal e federal.

Como desdobramentos futuros, destacam-se a implementação de modelos preditivos para estimativa
de indicadores em exercícios subsequentes, a geração automatizada de relatórios periódicos
em formato PDF e a extensão do chat analítico com suporte a memória de conversação para
análises de múltiplos turnos.

\bibliographystyle{elsarticle-harvard}
\bibliography{references}

\end{document}
